\documentclass[10pt,a4paper]{amsart}
\usepackage[margin=19mm]{geometry}
\usepackage{amsmath,amssymb,mathtools}
\usepackage[scr=rsfs,cal=euler]{mathalfa}
\usepackage{xfrac}
\usepackage[unicode,hidelinks,bookmarks=false]{hyperref}
\newcommand{\ie}{i.e.\ }
\newcommand{\cf}{cf.\ }
\newcommand{\resp}{respectively\ }
\newcommand{\Subsection}{\subsection}
\providecommand{\nobreakdash}{\nobreak\hbox{-}}
\setlength{\emergencystretch}{2em}
\multlinegap0pt
\usepackage{amssymb}
\newtheorem{lemma}{Lemma}
\title{Weighted Ces\`aro type operators on the weighted Bergman spaces in the unit ball}
\author{Jiaxin Pan, Cezhong Tong, Zicong Yang}
\date{Source: arXiv:2609.03243v1 (CC BY 4.0)}
\begin{document}
\maketitle

\noindent\emph{Source and licence.} Weighted Ces\`aro type operators on the weighted Bergman spaces in the unit ball. Version arXiv:2609.03243v1 (CC BY 4.0), distributed by arXiv under the Creative Commons Attribution 4.0 International licence, http://creativecommons.org/licenses/by/4.0/, as recorded in the arXiv licence metadata for this exact version and shown on the abstract page https://arxiv.org/abs/2609.03243v1. Full licence notice: \texttt{licenses/arxiv-2609.03243-CC-BY-4.0.txt}.\par\smallskip

\noindent\textit{Editorial scope.} Counted unit: the article's lemma lemma4.3 (source0.tex lines 1289--1298). The article's complete original proof of it is delivered with it (source0.tex lines 1300--1404, one \texttt{proof} environment of 105 lines). No mathematics is written, repaired or re-derived: the statement and the proof are the article's own lines, in the article's own notation, and no retained line is substituted or re-flowed.  Retained uncounted as the source-local context the proof needs: lines 423--432; lines 1209--1215; lines 1258--1263.  Nothing was omitted from any retained span, so the package carries no omission record.  No retained line contains a citation, so the wrapper carries no bibliography and no bibliography entry.  No retained line contains a figure environment, an image-inclusion command or drawing code, so no image is delivered and the package declares no asset dependency.  Editorial formatting only, in addition: an AMS \texttt{amsart} preamble that restates the article's own declarations and macros used by the retained lines, namely the environments \texttt{lemma} on separate counters, together with the standard packages those lines need.  The retained environments print in this document as Lemma 1 in order of appearance, whereas the article prints the same items with the numbering of its own full document.  The complete native source of the article as distributed in the arXiv e-print for this exact version is bundled unchanged as \texttt{sources/209280/source0.tex}.
\begin{lemma}  \label{lemma2.7}
  Let $\lambda\in\mathbb{R}$, $\theta>0$ and $\mu$ be a finite positive Borel measure on $[0,1)$. Then
\begin{equation}\label{equa2.1}
\int_{0}^1\Big(\frac{\hat{\mu}(r)}{(1-r)^{\lambda}}\Big)^{\theta}\frac{{\rm d}r}{1-r}<\infty
\end{equation}
if and only if 
\begin{equation}\label{equa2.2}
\sum_{k=1}^{\infty}\Big(\frac{\hat{\mu}(r_k)}{\delta_k^{\lambda}}\Big)^{\theta}<\infty.
\end{equation}
\end{lemma}

\begin{lemma}\label{lemma4.1}
Let $\alpha>-1$, $\theta>0$ and $\mu$ be a positive Borel measure on $[0,1)$. Then
\begin{equation}\label{equa4.1}
\int_{\mathbb{B}_n}\left|V_{\mu}^{\xi}(z)\right|^{\theta}{\rm d}v_{\alpha}(z)\simeq \int_{\mathbb{D}}\left|V_{\mu}(u)\right|^{\theta}(1-|u|^2)^{n+\alpha-1}{\rm d}A(u),
\end{equation}
where ${\rm d}A$ is the normalized area measure on $\mathbb{D}$.
\end{lemma}

\begin{lemma}\label{lemma4.2}
For $k\geq 1$, let $E_k=\{u\in\mathbb{D}: \delta_{k+1}<|1-u|\leq \delta_k\}$. Then 
\begin{equation*}
\int_{E_k}(1-|u|^2)^{n+\alpha-1}{\rm d}A(u)\simeq \delta_{k}^{n+1+\alpha}.
\end{equation*}
\end{lemma}

\begin{lemma}\label{lemma4.3}
Let $\alpha>-1$, $\beta>0$ and   $\theta>0$. Set
\[
    \lambda=\beta-\frac{n+1+\alpha}{\theta}.
\]
Suppose that $\mu$ is a finite positive Borel measure on $[0,1)$. Then $V_{\mu}^{\xi}\in L_{\alpha}^{\theta}(\mathbb{B}_n)$ if and only if
\begin{equation}\label{equa4.3}
\int_{0}^1\Big(\frac{\hat{\mu}(r)}{(1-r)^{\lambda}}\Big)^{\theta}\frac{{\rm d}r}{1-r}<\infty.
\end{equation}
\end{lemma}

\begin{proof}
By Lemma \ref{lemma2.7}, the integral condition in \eqref{equa4.3} is equivalent to
\[
    \sum_{k=1}^{\infty}\left(\frac{\hat{\mu}(r_k)}{\delta_k^{\lambda}}\right)^{\theta}<\infty,
\]
where $r_k=1-2^{-k}$, $\delta_k=2^{-k}$. Denote by $h_k=\hat{\mu}(r_k)/\delta_k^{\lambda}$. We only need to establish the equivalence between the condition   $\{h_k\}\in \ell^{\theta}$  and the condition $V_{\mu}^{\xi}\in L_{\alpha}^{\theta}(\mathbb{B}_n)$.

$\bullet$(i) Assume first that   $\{h_k\}\in \ell^{\theta}$. Let
\[
    V_{\mu,0}(u)=\int_{[0,\frac{1}{2})}\frac{{\rm d}\mu(t)}{|1-tu|^{\beta}},\quad V_{\mu,1}(u)=\int_{[\frac{1}{2}, 1)}\frac{{\rm d}\mu(t)}{|1-tu|^{\beta}},\quad u\in\mathbb{D}.
\]
By Lemma \ref{lemma4.1}, 
\begin{equation*}
\int_{\mathbb{B}_n}\left|V_{\mu}^{\xi}(z)\right|^{\theta}{\rm d}v_{\alpha}(z)\simeq \int_{\mathbb{D}}\left(|V_{\mu,0}(u)|^{\theta}+|V_{\mu,1}(u)|^{\theta}\right)(1-|u|^2)^{n+\alpha-1}{\rm d}A(u).
\end{equation*}
Clearly, 
\[
    \int_{\mathbb{D}}\left|V_{\mu,0}(u)\right|^{\theta}(1-|u|^2)^{n+\alpha-1}{\rm d}A(u)<\infty,
\]
since $\mu$ is finite and $n+\alpha-1>-1$. Hence we need to prove the finiteness of the integral $\int_{\mathbb{D}}|V_{\mu,1}(u)|^{\theta}(1-|u|^2)^{n+\alpha-1}{\rm d}A(u)$.

For $j\geq 1$, write $E_j=\{u\in\mathbb{D}: \delta_{j+1}<|1-u|\leq \delta_j\}$, $I_j=[r_j,r_{j+1})$ and $\mu_{j}=\mu(I_j)$. If $u\in E_k$ and $t\in I_j\, (k,j\geq 1)$, then $1-t\simeq \delta_j$, $|1-u|\simeq \delta_k$ and $t\geq 1/2$. Since ${\rm Re}(1-u)>0$, we get
\begin{align}\label{equa4.4}
|1-tu|^2=\left|1-t+t(1-u)\right|^2\geq (1-t)^2+t^2|1-u|^2\simeq \delta_j^2+\delta_k^2.
\end{align}
Hence, $|1-tu|\gtrsim \max\{\delta_j,\delta_k\}$. Therefore, when $u\in E_k$, we obtain
\begin{align*}
V_{\mu,1}(u)=\sum_{j=1}^{\infty}\int_{I_j}\frac{{\rm d}\mu(t)}{|1-tu|^{\beta}}\lesssim \sum_{j=1}^{\infty}\frac{\mu_j}{\max\{\delta_j,\delta_k\}^{\beta}}=\sum_{j<k}\frac{\mu_j}{\delta_j^{\beta}}+\sum_{j\geq k}\frac{\mu_j}{\delta_k^{\beta}}.
\end{align*}
Note that $\mu_j\leq \hat{\mu}(r_j)=h_j\cdot \delta_j^{\lambda}$, so we have
\begin{equation*}
\begin{split}
\sum_{j<k}\frac{\mu_j}{\delta_j^{\,\beta}}\leq \sum_{j<k}h_j\cdot \delta_j^{\,\lambda-\beta}=\sum_{j<k}h_j\cdot \delta_j^{\,-(n+1+\alpha)/\theta}=2^{k(n+1+\alpha)/\theta} \sum_{j<k}2^{(j-k)(n+1+\alpha)/\theta}\cdot h_j.
\end{split}
\end{equation*}
Besides,
\begin{equation*}
\sum_{j\geq k}\frac{\mu_j}{\delta_k^{\,\beta}}=\frac{\hat{\mu}(r_k)}{\delta_k^{\,\beta}}=h_k\cdot\delta_k^{\,\lambda-\beta}=2^{k(n+1+\alpha)/\theta}\cdot h_k.
\end{equation*}
Thus, when $u\in E_k$, we get
\begin{equation*}
V_{\mu,1}(u)\lesssim 2^{k(n+1+\alpha)/\theta}\cdot \widetilde{h}_k,
\end{equation*}
where $\widetilde{h}_k=\sum_{j\leq k}2^{(j-k)(n+1+\alpha)/\theta}\cdot h_j$. Then using Lemma \ref{lemma4.2}, we obtain
\begin{align*}
\int_{E_k}\left|V_{\mu,1}(u)\right|^{\theta}(1-|u|^2)^{n+\alpha-1}{\rm d}A(u)\lesssim \widetilde{h}_k^{\,\theta}.
\end{align*}

Let $h_0=0$ and $\mathbf{h}=\{h_k\}_{k\geq 0}$. Choose a sequence $\mathbf{a}=\{a_k\}$ such that
\[
     a_k=2^{-k(n+1+\alpha)/\theta},\quad k\geq 0.
\]
Then 
\[
    \widetilde{h}_k=\sum_{j=0}^{k}a_j\cdot h_{k-j}.
\]
  If  $\theta\geq 1$, the discrete Young inequality on $\ell^{\theta}$ yields
\[
    \left\|\{\widetilde{h}_k\}\right\|_{\ell^{\theta}}=\big\|\mathbf{a}*\mathbf{h}\big\|_{\ell^{\theta}}\leq \big\|\mathbf{a}\big\|_{\ell^1}\cdot \big\|\mathbf{h}\big\|_{\ell^{\theta}}\lesssim \big\|\mathbf{h}\big\|_{\ell^{\theta}}.
\]
 If $0<\theta<1$, subadditivity gives, with $A=n+1+\alpha>0$,
\[
    \widetilde h_k^{\,\theta}
    \leq \sum_{j=1}^{k}2^{-(k-j)A}h_j^{\,\theta}.
\]
Consequently,
\[
    \sum_{k=1}^{\infty}\widetilde h_k^{\,\theta}
    \leq \sum_{j=1}^{\infty}h_j^{\,\theta}
          \sum_{m=0}^{\infty}2^{-mA}
    \lesssim \sum_{j=1}^{\infty}h_j^{\,\theta}.
\]
 Consequently, 
\begin{equation}\label{equa4.5}
\begin{split}
&\quad \int_{\mathbb{D}\cap \big\{|1-u|<1/2\big\}}\left|V_{\mu,1}(u)\right|^{\theta}(1-|u|^2)^{n+\alpha-1}{\rm d}A(u)\\
&=\sum_{k=1}^{\infty}\int_{E_k}\left|V_{\mu,1}(u)\right|^{\theta}(1-|u|^2)^{n+\alpha-1}{\rm d}A(u)\lesssim \sum_{k=1}^{\infty}\widetilde{h}_k^{\,\theta}\lesssim \sum_{k=1}^{\infty}h_k^{\,\theta}<\infty.
\end{split}
\end{equation}
On the other hand, if $|1-u|\geq 1/2$ and $t\geq 1/2$, then
\[
    |1-tu|\gtrsim (1-t)+t|1-u|\geq \frac{1}{4}.
\]
It follows that 
\begin{align*}
&\quad\int_{\mathbb{D}\cap\big\{|1-u|\geq 1/2\big\}}\left|V_{\mu,1}(u)\right|^{\theta}(1-|u|^2)^{n+\alpha-1}{\rm d}A(u)\\
&\lesssim \left[4^{\beta}\mu([1/2,1))\right]^{\theta}\cdot\int_{\mathbb{D}}(1-|u|^2)^{n+\alpha-1}{\rm d}A(u)<\infty.
\end{align*}
This, together with \eqref{equa4.5}, shows that
\[
    \int_{\mathbb{D}}\left|V_{\mu,1}(u)\right|^{\theta}(1-|u|^2)^{n+\alpha-1}{\rm d}A(u)<\infty.
\]

$\bullet$(ii) Assume that $V_{\mu}^{\xi}\in L_{\alpha}^{\theta}(\mathbb{B}_n)$. If $u\in E_k$ and $t\in [r_k,1)$, then $|1-tu|\leq (1-t)+t|1-u|\lesssim \delta_k$. Hence
\[
    V_{\mu}(u)\geq \int_{[r_k,1)}\frac{{\rm d}\mu(t)}{|1-tu|^{\beta}}\gtrsim \frac{\hat{\mu}(r_k)}{\delta_k^{\,\beta}}=h_k\cdot \delta_k^{-(n+1+\alpha)/\theta}.
\]
Then by Lemma \ref{lemma4.1} and Lemma \ref{lemma4.2}, we obtain
\begin{align*}
\int_{\mathbb B_n}\left|V_{\mu}^{\xi}(z)\right|^{\theta}{\rm d}v_{\alpha}(z)&\simeq\int_{\mathbb{D}}\left|V_{\mu}(u)\right|^{\theta}(1-|u|^2)^{n+\alpha-1}{\rm d}A(u)\\
&\geq \sum_{k=1}^{\infty}\int_{E_k}\left|V_{\mu}(u)\right|^{\theta}(1-|u|^2)^{n+\alpha-1}{\rm d}A(u)\\
&\gtrsim \sum_{k=1}^{\infty}\delta_k^{\,n+1+\alpha}\cdot \big(h_k\cdot \delta_k^{-(n+1+\alpha)/\theta}\big)^{\theta}=\sum_{k=1}^{\infty}h_k^{\,\theta}.
\end{align*}
Consequently, $\{h_k\}\in \ell^{\,\theta}$. The proof is now complete.
\end{proof}

\end{document}
